\chapter{Access Control, Audit, and MCP in the Enterprise}
\label{ch:access-control}

In an agentic application the chain of responsibility runs user, model, host,
client, server, data. Every link is a place where accountability can go missing
if nobody wrote it down, and the components most likely to be blamed afterwards
are usually the ones with the least ability to record anything useful.

This chapter is about writing it down, and about the access controls that make
the record worth having. An audit log full of actions that should never have
been permitted is a precise description of an incident.

\section{Scoped authorization, in practice}

Chapter~\ref{ch:authorization} did the wire mechanics. Here is what to do with
them.

Authorization decisions land at three different granularities, and a system of
any size needs all three. Each one catches something the other two cannot see.

\subsection{Per-tool}

The coarsest and the easiest. A scope per tool, or per group of tools.

\begin{py}
REQUIRED_SCOPE = {
    "assess_account_risk": "risk:read",
    "rank_portfolio_risk": "risk:read",
    "underwrite_loan":     "risk:write",
    "review_transaction":  "compliance:write",
    "export_audit_log":    "audit:read",
}

class ScopedRiskServer(Server):
    """The risk server, with per-tool scopes and row-level account access."""

    def visible_tools(self, ctx: RequestContext) -> list[Tool]:
        """Override to filter by scope. The set may vary by authorization,
        never by connection."""
        scopes = set((ctx.auth or {}).get("scopes", []))
        return [t for t in super().visible_tools(ctx)
                if not REQUIRED_SCOPE.get(t.name)
                or REQUIRED_SCOPE[t.name] in scopes]
\end{py}

Two properties worth noting.

The filtering happens in \texttt{visible\_tools}, so an unauthorized tool is \emph{invisible}. The model never sees it,
never plans around it, and never wastes a turn discovering it cannot call it.

And \texttt{super()} sorts first, so the filtered list stays deterministic and
the prompt-cache property from Chapter~\ref{ch:tools} survives.

\begin{dangerbox}{Invisible is not the same as protected}
Filtering the catalogue is a usability feature. It is not access control.

A caller can invoke \texttt{tools/call} for a tool that was never listed. So the
handler must check authorization too. If your only check is in
\texttt{visible\_tools}, you have built a menu, not a lock.
\end{dangerbox}

\subsection{Per-resource}

\begin{py}
# sketch: the shipped handler with an authorization check added inline
@server.template("meridian://accounts/{account_id}/summary", ...)
def account_summary(ctx: RequestContext, uri: str, params: dict):
    account_id = params["account_id"]
    if not can_read_account(ctx.auth, account_id):
        # Same error as "does not exist". Distinguishing them tells an
        # attacker which account ids are real.
        raise ResourceNotFound(uri)
    ...
\end{py}

That comment is the point. Returning ``forbidden'' for accounts that exist and
``not found'' for those that do not turns your error responses into an
enumeration oracle. Return the same thing for both.

\subsection{Per-request, which is row-level security}

The finest grain, and the one multi-tenancy actually rests on: the same tool,
called by two principals with the same arguments, has to return different rows.

\begin{py}
# sketch: the shipped handler with row-level filtering added inline
def rank_portfolio_risk(ctx: RequestContext):
    ids = ctx.arguments.get("accountIds") or []
    # Intersect what was asked for with what this principal may see, before
    # touching the data. The model can request anything; the server decides.
    permitted = accounts_visible_to(ctx.auth)
    ids = [a for a in ids if a in permitted]
    ...
\end{py}

\keyidea{The model can ask for anything. It has read attacker-influenced text and
it is not an authorization component. Every access decision belongs on the server,
against the credential on the request, and never against an argument the model
supplied.}

\section{Context isolation between tenants}

The architecture gives you some of this. The rest you build.

\begin{table}[htbp]
\centering
\small
\begin{tabularx}{\textwidth}{@{}lXl@{}}
\toprule
\thead{Leak path} & \thead{Control} & \thead{Where} \\
\midrule
Shared cache entry & Fold auth context into every cache key & host \\
Shared task store & Authorize \texttt{tasks/get} per caller & server \\
MRTR state replay & Bind \texttt{requestState} to the principal & server \\
Handle guessing & Opaque handles, authorized on every use & server \\
Enumerable URIs & Non-sequential ids; identical not-found errors & server \\
Error message detail & Say less to callers, log more internally & server \\
\bottomrule
\end{tabularx}
\caption{Six paths. The first is the one that fails silently and at scale.}
\label{tab:isolation-paths}
\end{table}

Meridian's cache is deliberately conservative:

\begin{py}
def cache_key(server, method, params, auth_context):
    """Method plus the params that affect the answer, plus who is asking.

    `auth_context` is folded in unconditionally rather than only for `private`
    entries. Doing it the other way round means a single mislabelled response
    leaks across users, and the cost of being wrong is far higher than the cost
    of a few duplicate entries.
    """
\end{py}

That trades some gateway-level sharing of public resources for making a
mislabelling bug harmless. For a financial system it is obviously the right
trade. Make the choice deliberately.

\begin{py}
def test_private_entries_are_keyed_by_auth_context(self):
    cache = ResultCache()
    result = {"resultType": "complete",
              "contents": [{"uri": "u", "text": "alice"}],
              "ttlMs": 60000, "cacheScope": "private"}
    cache.put("s", "resources/read", {"uri": "u"}, result, auth_context="alice")
    self.assertIsNotNone(cache.get("s", "resources/read", {"uri": "u"}, "alice"))
    self.assertIsNone(cache.get("s", "resources/read", {"uri": "u"}, "bob"))
\end{py}

\section{Human in the loop as a control}

Chapter~\ref{ch:mrtr} covered the mechanism. Here it is as a security boundary,
and the distinction that makes it one.

Meridian requires a named officer for any flagged transaction. Not a checkbox. A
name, which goes in the audit log, attached to a decision.

\begin{py}
"officer": {"type": "string", "title": "Officer", "minLength": 2},
"decision": {
    "type": "string", "title": "Decision",
    "enum": ["clear", "escalate", "hold"],
    "default": "escalate",
},
"note": {"type": "string", "title": "Note", "maxLength": 500},
\end{py}

Four design decisions in that schema:

\textbf{The officer is named.} An approval without an approver is a log entry
nobody owns.

\textbf{The decision is an enum.} Three outcomes, no free text, so the audit log
is queryable.

\textbf{The default is \texttt{escalate}, not \texttt{clear}.} A user who clicks
through without reading escalates. Defaults are policy, and this one fails safe.

\textbf{There is a note field.} Because ``why'' is what the reviewer will want in
eighteen months, and the only moment anyone knows it is now.

\subsection{Where the gate belongs}

\begin{table}[htbp]
\centering
\small
\begin{tabularx}{\textwidth}{@{}lXX@{}}
\toprule
\thead{Gate at} & \thead{Good for} & \thead{Weakness} \\
\midrule
Host consent prompt & Anything with side effects & The user learns to click \\
Server MRTR elicitation & Domain rules the host cannot know & Costs a round trip \\
Both & High-value operations & Two prompts unless coordinated \\
\bottomrule
\end{tabularx}
\caption{Server-side gates are the stronger ones, because they hold regardless of
which host is calling.}
\label{tab:gate-placement}
\end{table}

That last point matters more than it looks. A host-side gate protects users of
\emph{your} host. A server-side gate protects everyone, including a user running a
host you have never seen, which in an MCP ecosystem is most of them.

\keyidea{Put approval gates for domain-critical operations on the server, not in
the host. The server is the only component that every caller shares.}

\section{Audit, designed for the regulator}

Most audit logs are debugger output with a compliance label. Here is what changes
when you write one for somebody who will read it under pressure, years later.

\begin{py}
def _audit(ctx: RequestContext, tool: str, subject: str, payload: dict) -> None:
    """Append-only, structured, and written for a regulator rather than a
    debugger."""
    AUDIT_LOG.append({
        "tool": tool,
        "subject": subject,
        "principal": (ctx.auth or {}).get("sub", "anonymous"),
        "client": ctx.client_info.name if ctx.client_info else "unknown",
        "protocolVersion": ctx.protocol_version,
        "traceparent": ctx.traceparent,
        "outcome": payload.get("outcome"),
        "officer": payload.get("officer"),
    })
\end{py}

\subsection{What a record must contain}

\begin{table}[htbp]
\centering
\small
\begin{tabularx}{\textwidth}{@{}lX@{}}
\toprule
\thead{Field} & \thead{Answers} \\
\midrule
\texttt{principal} & Who, authoritatively, from the token \\
\texttt{tool}, \texttt{subject} & What, and to which object \\
timestamp & When \\
\texttt{traceparent} & What else was part of the same request \\
\texttt{outcome} & What happened \\
\texttt{officer} & Who authorised it, when a human did \\
\texttt{client} & Which software, for display only \\
\texttt{protocolVersion} & Which contract was in force \\
\bottomrule
\end{tabularx}
\caption{Eight fields. \texttt{traceparent} is the one that turns a list of
events into a reconstructable story.}
\label{tab:audit-fields}
\end{table}

Two notes on the last two rows.

\texttt{client} is self-reported and unverified. Log it, because it is useful for
support. Never make a decision with it, and never present it in a report as
though it were established fact.

\texttt{protocolVersion} matters because behaviour changes between revisions. In
2028 somebody will ask why a 2026 request behaved a particular way, and the
answer will be ``because it was \texttt{2026-07-28}''.

\subsection{What Chapter~\ref{ch:threats} proved was missing}

The pen test caught the fan-out and could not explain it. The log recorded every
action and nothing about the \emph{input} that caused them.

\begin{py}
AUDIT_LOG.append({
    ...
    # What entered the model's context, and where it came from. Without this,
    # an injection is invisible in the record: you see the actions and not the
    # sentence that produced them.
    "inputDigest": hashlib.sha256(result_text.encode()).hexdigest()[:32],
    "inputSource": f"{server_label}.{tool_name}",
    "inputBytes": len(result_text),
})
\end{py}

Hashes rather than content, for a reason: the content may be private, and
retaining it turns your audit log into a second copy of the data you were
protecting. A hash plus a source plus a length is enough to prove that a
particular result preceded a particular behaviour, which is what an incident
review needs.

\subsection{Properties, not just fields}

\begin{itemize}[leftmargin=1.6em, itemsep=3pt]
\item \textbf{Append-only.} If it can be edited it is a database, not an audit
log. \secref{tamper-evident} is about making that true rather than declaring it.
\item \textbf{Written before the action.} Meridian audits before executing, so
the record exists even when the call times out.
\item \textbf{Independently queryable.} By principal, by subject, by trace, by
time.
\item \textbf{Retained past the incident.} Compromises are found late. A 30-day
retention is a 30-day investigation limit.
\item \textbf{Not the same store as your application data.} An attacker who
reaches your database should not thereby reach the record of what they did.
\end{itemize}

\subsection{Making ``append-only'' true}
\label{sec:tamper-evident}

That first bullet is the one everybody writes down and nobody implements. A
Python list is not append-only. A Postgres table is not append-only. Both are
writable by anybody who can write, which includes the attacker whose actions you
were recording, and editing the log is what a competent one does on the way out.

You cannot stop the edit at the storage layer without giving up the ability to
write at all. You can make it \emph{detectable}, which is enough, because an
attacker who cannot cover their tracks is an attacker whose intrusion gets
reconstructed.

The mechanism is a hash chain. Each entry carries a digest computed over its own
contents and the previous entry's digest:

\begin{py}
def digest(entry: dict, previous: str) -> str:
    """Hash one entry together with its predecessor's digest.

    Serialised with sorted keys, because two encoders that disagree on field
    order would produce different digests for identical entries and the chain
    would fail to verify for a reason that has nothing to do with tampering.
    """
    body = json.dumps(entry, sort_keys=True, separators=(",", ":"),
                      default=str)
    return hashlib.sha256(f"{previous}{body}".encode()).hexdigest()
\end{py}

Every entry now depends on every entry before it. Edit entry 40 and entry 41's
recorded \texttt{previous} no longer matches entry 40's recomputed digest, and
so on to the end. \texttt{verify} walks the chain and reports the first index
where it breaks, which is both an alarm and a pointer at what was touched:

\begin{py}
def test_editing_an_entry_is_detected(self):
    self.chain[2]["subject"] = "T-forged"
    ok, bad = self.chain.verify()
    self.assertFalse(ok)
    self.assertEqual(bad, 2)
\end{py}

Deletion and reordering are caught by the same property, which is worth noting
because they are the operations a naive integrity check misses. A per-row
checksum detects edits and says nothing about a row that is no longer there.

Now the honest limitation, because a control whose limits you have not stated is
a control somebody will over-trust. An attacker who can write to the log can
also recompute the whole chain from their edit onward. Verification then passes
against a log that lies.

What closes that is one number. Publish the head digest, which covers every
entry beneath it, somewhere the attacker does not control: a different account,
a log-shipping pipeline, a customer's system, an hourly line in a chat channel.
Rewriting the chain now requires also rewriting every published head, and the
first one they miss is the alarm.

\keyidea{Chaining is what makes publishing cheap. One digest an hour covers
every entry written in that hour, so tamper-evidence costs a cron job rather
than a replicated write path.}

\section{Wrapping legacy systems}

The practical enterprise problem. There is a mainframe. There is a SOAP service
somebody wrote in 2009. There is an API with 340 endpoints.

\subsection{Anti-corruption layer}

Do not expose the legacy model. Chapter~\ref{ch:tools} measured what happens: the
REST mirror cost 4,371 extra tokens per turn and \$26 per thousand tasks in pure
waste.

\begin{table}[htbp]
\centering
\small
\begin{tabularx}{\textwidth}{@{}XX@{}}
\toprule
\thead{Legacy shape} & \thead{MCP shape} \\
\midrule
\texttt{GET /accounts/\{id\}} plus four more calls & \texttt{assess\_account\_risk} \\
\texttt{POST /jobs} then poll \texttt{/jobs/\{id\}} & one tool returning a task \\
SOAP envelope with 40 optional fields & a schema with the six that matter \\
Error codes in a 200 response body & tool execution errors with actionable text \\
\bottomrule
\end{tabularx}
\caption{The translation is the work. A thin wrapper over a legacy API produces a
legacy-shaped MCP server, which is the worst of both.}
\label{tab:anti-corruption}
\end{table}

\subsection{Strangler fig}

Migrating from bespoke function calling, without a rewrite:

\begin{enumerate}[leftmargin=1.6em, itemsep=3pt]
\item Stand up an MCP server beside the existing integration. Both work.
\item Move one workflow. Measure it: round trips, tokens, cost, success rate.
\item Move the rest, one at a time, each with a measurement.
\item Delete the old integration when nothing calls it.
\end{enumerate}

Step two is the one to insist on. If the first migrated workflow is not measurably
better, the problem is your tool design and moving forty more will multiply it.

\section{Registries and change management}

At enterprise scale, a private registry, and the governance questions it answers:

\begin{table}[htbp]
\centering
\small
\begin{tabularx}{\textwidth}{@{}lX@{}}
\toprule
\thead{Question} & \thead{Answered by} \\
\midrule
Which servers exist? & The registry \\
Who owns this one? & A required ownership field \\
What data does it touch? & A data-classification field \\
Has the App template been reviewed? & An approval record with a hash \\
What changed last Tuesday? & A version history \\
Which hosts may connect? & An allowlist \\
\bottomrule
\end{tabularx}
\caption{A registry is a governance artefact that happens to also do discovery.}
\label{tab:registry}
\end{table}

Two policies worth adopting.

\textbf{Approve tool definitions, not just servers.} Pin the hash of what you
approved. Chapter~\ref{ch:threats} explained why: without pinning,
\texttt{listChanged} is a rug-pull mechanism.

\textbf{Give your own surface a deprecation window.} The specification promises
you twelve months. Your internal consumers deserve the same, and the discipline
is identical: mark it Deprecated, publish the migration, wait, then remove.

\section{Measuring success}

These are the numbers that survive contact with a steering committee, which is a
different requirement from the numbers that help you engineer the thing.

\begin{table}[htbp]
\centering
\small
\begin{tabularx}{\textwidth}{@{}lXl@{}}
\toprule
\thead{KPI} & \thead{Definition} & \thead{Source} \\
\midrule
Task success rate & Completed without human rescue & evals \\
p95 task latency & What impatient users experience & traces \\
Cost per task & Total model spend divided by tasks & accounting \\
Context coherence & Tasks completing within the token budget & loop stats \\
Integration cost & Days to connect a new system & project tracking \\
Consent friction & Prompts per task, and approval rate & host audit \\
\bottomrule
\end{tabularx}
\caption{The last two are the ones nobody tracks. Both predict whether the system
survives contact with real users.}
\label{tab:kpis}
\end{table}

Integration cost is the one that justifies MCP existing. Before, connecting a new
system meant an integration per host. After, it means one server. If that number
is not falling, the protocol is not paying for itself and somebody should ask
why.

Consent friction is the leading indicator of a security control that is about to
stop working. If users approve 99.8\,\% of prompts, they are not reading them,
and your human-in-the-loop is a human-shaped rubber stamp.

\section{Meridian, final state}

\begin{plain}
                        [ users ]
                            |
                   [ host: consent, audit,
                     namespacing, budgets ]
                            |
                   [ supervisor agent ]
                            |
        +-------------------+-------------------+
        |                   |                   |
  [ risk agent ]    [ compliance agent ]  [ fraud agent ]
        |                   |                   |
  [ risk server ]   [ compliance server ] [ fraud server ]
        |                   |                   |
        +-------------------+-------------------+
                            |
                   [ marketdata server ]
                     public, gateway-cached
\end{plain}

\begin{measurebox}{The complete measured system}
\begin{tabular}{@{}lrr@{}}
\toprule
& \thead{Chapter 1} & \thead{Chapter 20} \\
\midrule
Loop iterations       & 4 & 2 \\
Wall clock            & 1{,}753\,ms & 1{,}014\,ms \\
Tokens per task       & 22{,}846 & 2{,}692 \\
Cost per task         & \$0.0690 & \$0.0086 \\
Setup requests (20 tasks) & 164 & 8 \\
Cache hit rate        & 0\,\% & 96.2\,\% \\
\midrule
Tests                 & \multicolumn{2}{r}{210, all passing} \\
Third-party runtime dependencies & \multicolumn{2}{r}{0} \\
\bottomrule
\end{tabular}

\vspace{6pt}
\footnotesize
Regenerate every row with \texttt{make bench} and \texttt{make test}.
\end{measurebox}

Every number is reproducible on your own machine, and every one of them came from a technique in this book.

\section{What to remember}

\begin{itemize}[leftmargin=1.6em, itemsep=3pt]
\item Filter the catalogue by scope for usability, and check authorization in the
handler for safety. The first is a menu, the second is a lock.
\item Identical errors for forbidden and nonexistent, or your errors enumerate
your data.
\item Every access decision uses the credential on the request, never an argument
the model supplied.
\item Fold the auth context into every cache key, public or private.
\item Put approval gates on the server. It is the only component every caller
shares.
\item Defaults are policy. Default to \texttt{escalate}.
\item Audit for the regulator: principal, subject, trace, outcome, approver, and
a digest of what entered the model's context.
\item Hash inputs rather than storing them, or your audit log becomes a second
copy of the data.
\item Wrap legacy systems with a translation, never a mirror.
\item Track integration cost and consent friction. Nobody else will.
\end{itemize}

That is the book. What remains is where it goes next.
